\documentclass[a4paper,14pt]{extreport}

% Русский язык и кодировка
\usepackage[T2A]{fontenc}
\usepackage[utf8]{inputenc}
\usepackage[english,russian]{babel}

% Математика и рисунки
\usepackage{amsmath, amssymb}
\usepackage{graphicx}
\usepackage{float}
\usepackage{geometry}
\usepackage{indentfirst}
\usepackage{hyperref}
\usepackage{enumitem}
\usepackage{caption}
\usepackage{booktabs}
\usepackage{xcolor}

% Настройка полей (левое не менее 2 см)
\geometry{
    a4paper,
    left=30mm,
    right=15mm,
    top=20mm,
    bottom=20mm,
}

% Настройка шрифтов и абзацев
\usepackage{setspace}
\setstretch{1.3} % межстрочный интервал 1,3-1,5
\setlength{\parindent}{1cm} % красная строка 1 см
\setlength{\parskip}{0pt}

% Настройка заголовков (без сокращений в названиях)
\usepackage{titlesec}
\titleformat{\chapter}
    {\normalfont\bfseries\centering}{\thechapter}{1em}{}
\titlespacing*{\chapter}{0pt}{-20pt}{20pt}

% Настройка подписей рисунков и таблиц по ГОСТу
\usepackage{caption}
\captionsetup[figure]{labelsep=period, justification=centering}
\captionsetup[table]{labelsep=period, justification=raggedright, singlelinecheck=off}

% Нумерация страниц
\usepackage{fancyhdr}
\pagestyle{fancy}
\fancyhf{}
\fancyfoot[C]{\thepage}
\renewcommand{\headrulewidth}{0pt}
\renewcommand{\footrulewidth}{0pt}

% Чтобы титульник и аннотации не нумеровались
\usepackage{chngcntr}
\usepackage{etoolbox}

% Настройка ссылок (ГОСТ)
\usepackage[backend=biber, style=gost-numeric, sorting=none, language=auto, autolang=other]{biblatex}
\addbibresource{references.bib}

% Адресация для гиперссылок
\hypersetup{
    colorlinks=true,
    linkcolor=black,
    citecolor=black,
    urlcolor=blue,
}

\begin{document}

% ------------------- ТИТУЛЬНЫЙ ЛИСТ -------------------
\begin{titlepage}
    \begin{center}
        {\small \textbf{Федеральное государственное автономное образовательное учреждение высшего образования «Национальный исследовательский университет «Высшая школа экономики}}\\[10mm]
        {\small \textbf{Факультет компьютерных наук}}\\
        {\small \textbf{Магистерская программа «Машинное обучение в цифровом продукте»}}\\[30mm]
        
        {\bfseries \large \textbf{КУРСОВАЯ РАБОТА}}\\[20mm]
    \end{center}
    
    \begin{flushleft}
        {На тему (рус.): \underline{Интерпретируемость больших языковых моделей}}\\[3mm]
        {На тему (англ.): \underline{Interpretability of Large Language Models}}\\[30mm]
    \end{flushleft}

    \begin{center}
        \hfill\begin{minipage}{0.5\textwidth}
            \raggedright
            Студент: \hrulefill\\
            (Ф.И.О., подпись)\\[5mm]
            Руководитель КР: \hrulefill\\
            (должность, звание, Ф.И.О., подпись)
        \end{minipage}\\[30mm]
        
        {\large Москва, 2026}
    \end{center}
\end{titlepage}
\newpage

% ------------------- АННОТАЦИИ (БЕЗ НОМЕРОВ) -------------------
\pagestyle{empty} % отключаем нумерацию
\section*{Аннотация}
\addcontentsline{toc}{chapter}{Аннотация}
В данной курсовой работе представлен подробный обзор методов механистической интерпретируемости нейронных сетей, основанный на трёх ключевых публикациях. Рассматриваются: математический фреймворк трансформерных цепей, включая QK/OV-декомпозицию и концепцию «индукционных голов»; метод разложения активаций языковых моделей с помощью разреженных автокодировщиков для получения моносемантических признаков; а также подход трассировки цепей на основе кросс-слойных транскодеров и графов атрибуции. Для каждого метода приводятся основная идея, ключевые формулы, архитектура и результаты применения. Работа предназначена для исследователей в области безопасности и интерпретируемости больших языковых моделей.


\vspace{15mm}

\section*{Abstract}
\addcontentsline{toc}{chapter}{Abstract}
This paper provides a detailed survey of mechanistic interpretability methods for neural networks, drawing on three foundational publications. We review: the mathematical framework for transformer circuits, including QK/OV circuit decomposition and the concept of induction heads; the sparse autoencoder approach for decomposing language model activations into monosemantic features; and the circuit tracing methodology based on cross-layer transcoders and attribution graphs. For each method, we present the core idea, key mathematical formulations, architectural details, and empirical results. The work is intended for researchers in the field of safety and interpretability of large language models.


\newpage

% ------------------- СОДЕРЖАНИЕ -------------------
\pagestyle{plain}
\tableofcontents
\newpage

% ------------------- ВВЕДЕНИЕ -------------------
\chapter*{Введение}
\addcontentsline{toc}{chapter}{Введение}
\pagestyle{plain}

Языковые модели, основанные на архитектуре трансформера~\cite{vaswani2017attention}, достигли выдающихся результатов во многих задачах обработки естественного языка. Такие модели, как GPT, Qwen,Deepseek и им подобные, способны генерировать связный текст, решать задачи и писать код. Вместе с тем их внутренние механизмы остаются в значительной мере непрозрачными: модель представляет собой миллиарды числовых параметров, не поддающихся непосредственному человеческому анализу.

Механистическая интерпретируемость ставит перед собой цель «реверс инжениринга» нейронных сетей, то есть понимания конкретных вычислительных механизмов, лежащих в основе наблюдаемого поведения. Это направление приобретает особую важность в контексте безопасности ИИ: систематическое понимание того, как именно модель принимает решения, позволяет обнаруживать нежелательное поведение, проверять соответствие целям разработчиков и предсказывать возможные сбои.

В данной работе рассматриваются взаимосвязанные направления механистической интерпретируемости, представленных в публикациях исследовательской группы Anthropic:


\newpage

% ------------------- ГЛАВА 1 -------------------
\chapter{Обзор методов интерпретируемости}

\section{Разреженные автокодировщики для поиска моносемантических признаков}

Данный раздел посвящён методу, предложенному в работе Bricken et al.~\cite{bricken2023monosemanticity}
«Towards Monosemanticity: Decomposing Language Models With Dictionary Learning» (Anthropic, 2023).
Метод направлен на решение фундаментальной проблемы механистической интерпретируемости~---
полисемантичности нейронов языковых моделей~--- с помощью разреженных автокодировщиков (Sparse Autoencoder, SAE).

\subsection{Мотивация: проблема полисемантичности}

Центральная проблема механистической интерпретируемости состоит в том, что нейроны
нейронных сетей, как правило, являются \textbf{полисемантическими}: один нейрон активируется
на несколько семантически несвязанных концепций. Авторы приводят наглядный пример из
однослойного трансформера: нейрон 2\,049 MLP-слоя сильно активируется на токены,
относящиеся к академическим цитатам, переводам, фрагментам прозы в кавычках,
обозначениям времени в японском тексте и последовательностям ДНК~--- совершенно разнородный набор концепций.

Такая полисемантичность существенно затрудняет интерпретацию: даже если исследователь
изучает активацию отдельного нейрона, он не может установить, какой именно концепт этот
нейрон «видит» в конкретном контексте.

Ключевой гипотезой, объясняющей это явление, служит \textbf{гипотеза суперпозиции}.

\subsection{Гипотеза суперпозиции}

Гипотеза суперпозиции~\cite{elhage2022superposition} утверждает, что нейронные сети
представляют \textit{больше признаков}, чем у них имеется нейронов, кодируя эти признаки
через линейные комбинации нейронов. Это возможно, если выполняются два условия:

\begin{itemize}
    \item \textbf{Разреженность признаков}: каждый признак активен лишь на малой доле
          входных данных, поэтому вероятность одновременной активации двух «соседних»
          признаков невелика;
    \item \textbf{Приблизительная ортогональность}: признаки кодируются в направлениях,
          которые примерно ортогональны друг другу, что минимизирует взаимную интерференцию.
\end{itemize}

Игрушечная модель суперпозиции~\cite{elhage2022superposition} демонстрирует: при
обучении линейной сети с $n$ нейронами представлять $m > n$ разреженных признаков
сеть «упаковывает» все признаки в суперпозицию, жертвуя точностью, но выигрывая в ёмкости.
Форма хранения зависит от соотношения $m/n$ и уровня разреженности: при высокой
разреженности признаки образуют примерно правильные многогранники в $n$-мерном пространстве.

Таким образом, полисемантичность нейронов является \textit{не случайным артефактом
обучения}, а предсказуемым следствием суперпозиции: нейрон одновременно участвует в
представлении нескольких признаков.

\subsection{Постановка задачи}

Авторы исследуют \textbf{однослойный трансформер} с 512-мерным пространством
активаций ($d = 512$), обученный на корпусе The Pile~\cite{gao2020pile}.
Объектом анализа служат \textbf{post-ReLU активации MLP-слоя}: вектор $\mathbf{x} \in \mathbb{R}^n$,
где $n = 512$~--- число нейронов в расширенном пространстве MLP ($d_{\mathrm{MLP}} = 4d = 2048$,
но авторы обозначают анализируемое измерение через $n$).

Задача формулируется следующим образом: найти \textbf{словарь признаков} из $m \gg n$
элементов~--- набор направлений $\{\mathbf{d}_i\}_{i=1}^{m}$ в пространстве $\mathbb{R}^n$~---
такой, что каждый вектор активаций $\mathbf{x}$ может быть приближённо представлен через
небольшое число активных признаков:
\begin{equation}
    \mathbf{x} \approx \mathbf{b}_{\mathrm{dec}} + \sum_{i=1}^{m} f_i(\mathbf{x}) \cdot \mathbf{d}_i,
    \label{eq:decomp}
\end{equation}
где $f_i(\mathbf{x}) \geq 0$~--- скалярное значение активации $i$-го признака, $\mathbf{b}_{\mathrm{dec}}$~--- смещение,
и при этом большинство $f_i(\mathbf{x})$ равны нулю (разреженность).

\subsection{Признаки как декомпозиция активаций}

Прежде чем перейти к архитектуре SAE, авторы формализуют понятие \textbf{хорошей декомпозиции}.
Согласно статье, декомпозиция активаций является качественной, если:

\begin{enumerate}
    \item \textbf{Точность}: восстановленные активации $\hat{\mathbf{x}}$ близки к исходным $\mathbf{x}$
          (малая ошибка реконструкции);
    \item \textbf{Разреженность}: среднее число ненулевых признаков на токен ($L_0$-норма)
          мало~--- в идеале от 1 до нескольких десятков;
    \item \textbf{Интерпретируемость}: каждый найденный признак активируется
          на семантически однородном множестве токенов и поддаётся человеческому объяснению;
    \item \textbf{Универсальность}: признаки, найденные независимыми запусками обучения,
          воспроизводят друг друга (т.\,е. не являются артефактами конкретной инициализации).
\end{enumerate}

Авторы обосновывают, почему нельзя использовать отдельные нейроны как «признаки»:
полисемантичность означает, что направление нейрона не соответствует ни одной
интерпретируемой концепции. PCA и ICA также не подходят: главные компоненты не
разреженны и не соответствуют семантическим концептам. Кластеризация не применима,
поскольку один токен может активировать несколько признаков одновременно. Метод NMF
требует неотрицательности исходных данных. Оптимальным решением оказывается
\textbf{разреженное кодирование} на переполненном (overcomplete) словаре.

\subsection{Архитектура разреженного автокодировщика}

SAE состоит из энкодера и декодера с общим промежуточным пространством размерности $m$.

\paragraph{Энкодер.}
Энкодер принимает вектор активаций $\mathbf{x} \in \mathbb{R}^n$, вычитает декодерное смещение
и применяет аффинное преобразование с функцией активации ReLU:
\begin{equation}
    \mathbf{f}(\mathbf{x}) = \mathrm{ReLU}\!\left(\mathbf{W}_{\mathrm{enc}}\,(\mathbf{x} - \mathbf{b}_{\mathrm{dec}}) + \mathbf{b}_{\mathrm{enc}}\right),
    \label{eq:encoder}
\end{equation}
где $\mathbf{W}_{\mathrm{enc}} \in \mathbb{R}^{m \times n}$~--- матрица весов энкодера,
$\mathbf{b}_{\mathrm{enc}} \in \mathbb{R}^{m}$~--- смещение энкодера,
$\mathbf{b}_{\mathrm{dec}} \in \mathbb{R}^{n}$~--- смещение декодера (вычитается до кодирования),
а $\mathbf{f}(\mathbf{x}) \in \mathbb{R}^{m}$~--- вектор активаций признаков.

\paragraph{Декодер.}
Декодер линейно восстанавливает исходный вектор из активаций признаков:
\begin{equation}
    \hat{\mathbf{x}} = \mathbf{W}_{\mathrm{dec}}\,\mathbf{f}(\mathbf{x}) + \mathbf{b}_{\mathrm{dec}},
    \label{eq:decoder}
\end{equation}
где $\mathbf{W}_{\mathrm{dec}} \in \mathbb{R}^{n \times m}$. Столбцы матрицы декодера обозначаются
$\mathbf{d}_i \in \mathbb{R}^{n}$ и называются \textbf{атомами словаря} (dictionary atoms):
каждый $\mathbf{d}_i$ задаёт направление $i$-го признака в пространстве активаций.

\paragraph{Ограничение на норму.}
Принципиальным архитектурным решением является требование \textbf{единичной нормы} столбцов декодера:
\begin{equation}
    \|\mathbf{d}_i\|_2 = 1 \quad \forall\, i \in \{1, \ldots, m\}.
    \label{eq:norm_constraint}
\end{equation}
Это ограничение предотвращает вырожденное решение, при котором декодерные веса
уменьшаются до нуля, а энкодер компенсирует это произвольно большими активациями.
Нормировка обеспечивает, что каждый атом словаря является \textit{единичным направлением},
а масштабная информация полностью сосредоточена в значениях $f_i(\mathbf{x})$.

\paragraph{Дизайнерские решения.}

\begin{itemize}
    \item \textbf{Pre-bias subtraction (вычитание смещения до кодирования).}
          Декодерное смещение $\mathbf{b}_{\mathrm{dec}}$ аппроксимирует среднее значение
          вектора активаций по всему датасету. Вычитая его перед кодированием, авторы
          центрируют данные относительно их среднего, позволяя энкодеру фокусироваться
          на отклонениях от среднего~--- именно они несут структурированную семантическую информацию.

    \item \textbf{Отсутствие разделения весов (no weight tying).}
          В отличие от классических автокодировщиков, где $\mathbf{W}_{\mathrm{enc}} = \mathbf{W}_{\mathrm{dec}}^\top$,
          в SAE $\mathbf{W}_{\mathrm{enc}}$ и $\mathbf{W}_{\mathrm{dec}}$ полностью независимы.
          Это позволяет энкодеру и декодеру использовать различные геометрические
          представления одного и того же пространства признаков.

    \item \textbf{ReLU как механизм разреженности.}
          Функция ReLU обнуляет отрицательные активации, обеспечивая неотрицательность
          $\mathbf{f}(\mathbf{x})$ и вводя неявную разреженность: лишь те признаки, у которых
          предактивационное значение положительно, участвуют в реконструкции.
\end{itemize}

\subsection{Функция потерь}

SAE обучается минимизировать следующую функцию потерь:
\begin{equation}
    \mathcal{L}(\mathbf{x}) = \underbrace{\|\mathbf{x} - \hat{\mathbf{x}}\|_2^2}_{\text{ошибка реконструкции}} + \lambda\,\underbrace{\|\mathbf{f}(\mathbf{x})\|_1}_{\text{штраф за плотность}},
    \label{eq:loss}
\end{equation}
где $\lambda > 0$~--- коэффициент регуляризации, управляющий компромиссом между
точностью реконструкции и разреженностью активаций признаков.

\paragraph{Ошибка реконструкции ($L_2$).}
Первое слагаемое принуждает SAE точно восстанавливать исходные активации. Если бы оно
было единственным, модель свелась бы к полному представлению $\mathbf{x}$ через $\mathbf{f}(\mathbf{x})$
без каких-либо ограничений разреженности.

\paragraph{$L_1$-штраф.}
Второе слагаемое штрафует за сумму абсолютных значений активаций признаков, поощряя
решения, в которых большинство $f_i(\mathbf{x})$ равны нулю. Норма $L_1$ является
\textbf{выпуклой релаксацией} нормы $L_0$ (числа ненулевых компонент) и при этом
остаётся дифференцируемой, что допускает обучение градиентными методами.

\paragraph{Роль гиперпараметра $\lambda$.}
\begin{itemize}
    \item При малом $\lambda$: SAE точно восстанавливает $\mathbf{x}$, но признаки плотные
          и плохо интерпретируемые;
    \item При большом $\lambda$: признаки разреженные, но реконструкция неточная
          (теряется часть информации об активациях);
    \item Оптимальное $\lambda$ находится эмпирически, ориентируясь на метрики
          интерпретируемости и качества реконструкции.
\end{itemize}

Полная функция потерь, усреднённая по батчу токенов:
\begin{equation}
    \mathcal{L} = \mathbb{E}_{\mathbf{x}}\!\left[\|\mathbf{x} - \mathbf{W}_{\mathrm{dec}}\,\mathbf{f}(\mathbf{x}) - \mathbf{b}_{\mathrm{dec}}\|_2^2 + \lambda \sum_{i=1}^{m} f_i(\mathbf{x})\right].
    \label{eq:full_loss}
\end{equation}

\subsection{Процедура обучения}

\paragraph{Данные.}
SAE обучается на активациях, собранных заранее при прогоне однослойного трансформера на
случайных токенах из корпуса The Pile~\cite{gao2020pile}~--- открытого датасета объёмом
$\sim$800\,ГБ, охватывающего широкий спектр текстовых жанров и языков. Обучение SAE
ведётся отдельно от обучения языковой модели~--- на зафиксированных активациях.

\paragraph{Размеры словарей.}
Авторы обучают семейство SAE с различными коэффициентами расширения $m/n$:
$1\times$, $2\times$, $4\times$, $8\times$, $16\times$, $32\times$ относительно $n = 512$.
Это даёт словари от 512 до $16\,384$ признаков. Основные результаты получены на
словаре $m = 4096$ ($8\times$).

\paragraph{Оптимизация.}
Используется оптимизатор Adam. Нормировка столбцов декодера поддерживается на каждом
шаге через проецирование градиента: после шага оптимизатора все столбцы $\mathbf{d}_i$
нормируются на единичную длину, что реализует ограничение~\eqref{eq:norm_constraint}.

\paragraph{«Мёртвые» признаки.}
В процессе обучения часть признаков может так и не начать активироваться («мёртвые» признаки,
dead features). Это происходит из-за того, что градиент через неактивные нейроны ReLU равен нулю.
Авторы используют эвристику периодической переинициализации мёртвых признаков, которая
помогает задействовать неиспользуемые атомы словаря.

\subsection{Метрики оценки}

Авторы предлагают комплексную систему оценки качества обученного SAE.

\subsubsection{Качество реконструкции}

\begin{itemize}
    \item \textbf{$L_2$-ошибка}: $\mathbb{E}\!\left[\|\mathbf{x} - \hat{\mathbf{x}}\|_2^2\right]$~---
          среднеквадратическое отклонение восстановленных активаций от исходных.
    \item \textbf{Доля объяснённых потерь} (fraction of loss recovered): насколько
          использование $\hat{\mathbf{x}}$ вместо $\mathbf{x}$ изменяет cross-entropy потери
          языковой модели. Если подставить реконструкции SAE обратно в модель вместо
          исходных активаций MLP, деградация потерь показывает, сколько информации
          утеряно при кодировании.
\end{itemize}

\subsubsection{Разреженность}

\begin{itemize}
    \item \textbf{$L_0$-норма}: среднее число ненулевых активаций $f_i(\mathbf{x})$ на токен.
          Является прямой мерой разреженности: чем меньше, тем более сжатым является
          представление. Для словаря $4\,096$ признаков авторы достигают $L_0 \approx 25$~---
          в среднем на токен активируется лишь 25 из 4096 признаков ($\approx 0.6\%$).
    \item \textbf{$L_1$-норма}: $\mathbb{E}\!\left[\|\mathbf{f}(\mathbf{x})\|_1\right]$~---
          используется в функции потерь, менее информативна для интерпретации.
\end{itemize}

\subsubsection{Интерпретируемость}

Это наиболее важная и нетривиальная метрика, для оценки которой авторы разработали
двухуровневый подход.

\paragraph{Ручная оценка.}
Исследователи изучают каждый признак вручную:
\begin{enumerate}
    \item Отбираются токены с наибольшей активацией $f_i(\mathbf{x})$ (топ-10--20 по значению);
    \item Выводится контекст каждого токена (окно в 20--30 символов);
    \item Исследователь формулирует гипотезу о том, что именно активирует признак;
    \item Гипотеза проверяется на случайных токенах с высокой активацией.
\end{enumerate}

\paragraph{Автоматическая интерпретируемость.}
Авторы применяют языковую модель (GPT-4 или Claude) для автоматизации процесса:
\begin{enumerate}
    \item Для данного признака $i$ отбираются $k$ токенов с наибольшей активацией
          (положительные примеры) и $k$ токенов с нулевой активацией (отрицательные примеры);
    \item Языковая модель генерирует \textit{объяснение}~--- гипотезу о семантике признака;
    \item На отдельной выборке из $k'$ токенов (удерживаемой выборке) языковую модель
          просят предсказать, будет ли признак активен для каждого токена;
    \item Вычисляются \textbf{detection score} (доля правильно угаданных активаций)
          и \textbf{false positive rate} (доля ложноположительных срабатываний).
\end{enumerate}

Чем выше detection score при низком false positive rate, тем более точным и
специфичным является объяснение признака.

\subsection{Результаты: найденные признаки}

Словарь из $4\,096$ признаков ($8\times$-расширение) демонстрирует богатое разнообразие
интерпретируемых концептов. Авторы описывают подробно несколько категорий.

\paragraph{Базовые числовые признаки.}
\begin{itemize}
    \item \textit{Base-10 digits}: активируется на цифры $0$--$9$ вне зависимости
          от контекста;
    \item \textit{Arabic numerals in non-Latin scripts}: арабские числа, встречающиеся
          внутри текстов на нелатинских системах письма (например, в персидских или
          арабских текстах);
    \item \textit{Numbers in decimal fractions}: числа, выступающие в роли дробной части
          десятичной записи.
\end{itemize}

\paragraph{Биологические и научные признаки.}
\begin{itemize}
    \item \textit{DNA sequences}: активируется на токены внутри последовательностей
          нуклеотидов (ATCG...); является одним из наиболее часто цитируемых примеров
          в статье как демонстрация специфичности найденных признаков;
    \item \textit{Periodic table elements}: символы химических элементов;
    \item \textit{Mathematical notation}: математические символы и обозначения.
\end{itemize}

\paragraph{Лингвистические признаки.}
\begin{itemize}
    \item \textit{Hebrew text}: токены в тексте на иврите;
    \item \textit{Arabic script}: арабское письмо;
    \item \textit{Cyrillic text}: кириллица;
    \item \textit{Legal language}: юридические термины и формулировки.
\end{itemize}

\paragraph{Программный код.}
\begin{itemize}
    \item \textit{Function arguments}: имена аргументов функций в коде;
    \item \textit{Variable assignment}: оператор присваивания в различных языках
          программирования;
    \item \textit{Python import statements}: инструкции импорта в Python.
\end{itemize}

\paragraph{Академический и публикационный стиль.}
\begin{itemize}
    \item \textit{Academic citations}: токены в академических ссылках (например, «et al.»,
          «ibid.», номера страниц в скобках);
    \item \textit{Bible quotations}: фрагменты библейских текстов;
    \item \textit{Prose quotations}: цитаты в художественных текстах.
\end{itemize}

Авторы подчёркивают: признаки, выявленные SAE, \textit{значительно более моносемантичны},
чем исходные нейроны. Большинство топ-активирующих токенов для каждого признака
относятся к семантически однородной категории, и эксперты-люди без труда формулируют
объяснение.

\subsection{Глобальный анализ}

Помимо анализа отдельных признаков, авторы исследуют глобальные свойства выученного словаря.

\subsubsection{Распределение частот активации}

Частота активации признаков ($\rho_i = P(f_i(\mathbf{x}) > 0)$) распределена приближённо
по степенному закону: большинство признаков редкие ($\rho_i < 0.1\%$),
и лишь немногие активируются на значительной доле токенов. Такое распределение
согласуется с предположением о разреженности: признаки соответствуют специфическим
концептам, которые встречаются в тексте нечасто.

\subsubsection{Дробление признаков (Feature Splitting)}

При увеличении размера словаря от $1\times$ до $32\times$ наблюдается \textbf{дробление признаков}:
признак, представлявший в малом словаре широкую концепцию, в большем словаре
разбивается на несколько более специфических подпризнаков. Например:
\begin{itemize}
    \item Словарь $512$: один признак «числа»;
    \item Словарь $4\,096$: признаки «арабские числа», «числа в скобках»,
          «числа в десятичных дробях», «числа в временны́х обозначениях» и т.\,д.
\end{itemize}
Дробление свидетельствует о том, что SAE выявляет иерархическую структуру признакового
пространства и что большие словари действительно «открывают» новые, более тонкие различия.

\subsubsection{Универсальность признаков (Feature Universality)}

Сравнение признаков, выученных в разных запусках SAE (с различными инициализациями и
разными коэффициентами расширения), показывает высокую степень \textbf{воспроизводимости}.
Признаки, найденные в одном словаре, имеют хорошо выраженный «двойник» в другом словаре
(высокая косинусная близость атомов $\mathbf{d}_i$). Это говорит о том, что признаки
отражают реальную структуру активационного пространства, а не являются артефактами
конкретного обучения.

\subsubsection{«Мёртвые» признаки}

Часть атомов словаря остаётся «мёртвой» (dead features)~--- никогда не активируется в
течение оценочного периода. При $8\times$-расширении доля мёртвых признаков составляет
порядка 20--30\%. Авторы интерпретируют это как ограничение текущего метода обучения,
а не как свидетельство того, что в активациях нет дополнительных структур.

\subsection{Феноменология признаков}

Авторы проводят детальный анализ нескольких признаков, демонстрируя типичный паттерн.
Для каждого признака строится:
\begin{itemize}
    \item \textbf{Гистограмма активаций}: крайне правостороннее распределение~---
          большинство токенов дают нулевую активацию, и лишь единичные токены показывают
          высокие значения;
    \item \textbf{Scatter-plot}: активации на осях двух признаков (показывает слабую
          корреляцию между признаками~--- признак независимости);
    \item \textbf{Список топ-токенов}: контексты с наибольшей активацией,
          позволяющие оценить семантику признака;
    \item \textbf{Список анти-токенов}: контексты с наиболее отрицательным
          предактивационным значением (до ReLU)~--- то, что «противоположно» признаку.
\end{itemize}

Анализ подтверждает: признаки с высоким detection score имеют чёткую семантику,
однородные топ-токены и статистически значимую отделимость от фона.

\subsection{Каузальное вмешательство в активации признаков}

Одним из ключевых вопросов является: \textit{являются ли найденные признаки реальными
вычислительными единицами, влияющими на поведение модели?} Корреляция признака с
токенами ещё не означает каузальности.

Для ответа авторы проводят \textbf{интервенционные эксперименты} (feature steering):
\begin{enumerate}
    \item Выбирается признак $i$ с известной семантикой;
    \item При прогоне токена $t$ активация признака искусственно фиксируется:
          $f_i(\mathbf{x}) \leftarrow c$ (где $c$ значительно превышает естественный уровень);
    \item Наблюдается изменение распределения следующего токена (вывода модели).
\end{enumerate}

Результаты показывают, что принудительная активация признака вызывает \textit{предсказуемое
и семантически согласованное} изменение вывода. Например, активация признака
«DNA sequences» побуждает модель генерировать нуклеотидные последовательности.
Это подтверждает, что найденные признаки~--- не просто корреляты, а \textbf{каузально
значимые вычислительные единицы} языковой модели.

\subsection{Интерфейс для анализа признаков}

Авторы создают \textbf{интерактивный веб-интерфейс} для просмотра обученных признаков.
Интерфейс позволяет:
\begin{itemize}
    \item Просматривать топ-токены для каждого признака;
    \item Искать признаки по ключевым словам (например, ввести «DNA» и найти соответствующий признак);
    \item Изучать корреляции между признаками;
    \item Видеть предсказания автоматической интерпретируемости.
\end{itemize}
Доступность такого инструмента существенно упрощает дальнейшие исследования
и проверку интерпретируемости.

\subsection{Обсуждение и ограничения}

\subsubsection{Масштабируемость}

Основным ограничением работы является масштаб: SAE обучается на \textit{однослойном}
трансформере с 512-мерными активациями. Применение метода к современным многослойным
моделям (например, GPT-4 или Claude с тысячами измерений) требует значительно большего
вычислительного ресурса и может потребовать новых архитектурных решений.

\subsubsection{«Является ли контекст признаком?»}

Авторы обсуждают философский вопрос: некоторые «признаки» кодируют не столько семантику
\textit{токена}, сколько \textit{контекст}. Например, признак может активироваться не
на сам токен «ATCG», а на любой токен в тексте, содержащем последовательность ДНК.
Это разграничение важно для понимания того, что именно представляют признаки~---
локальную или глобальную информацию.

\subsubsection{Дополнительные свойства.}
\begin{itemize}
    \item \textbf{Признаки сами могут быть полисемантичными}: в малых словарях один признак
          объединяет семантически различные, но статистически коррелирующие концепты.
          Дробление при увеличении словаря частично решает эту проблему.
    \item \textbf{Влияние на кросс-энтропию}: подстановка реконструкций SAE вместо
          исходных активаций увеличивает кросс-энтропийные потери языковой модели,
          что указывает на неизбежные потери информации при сжатии. Авторы показывают,
          что большие словари дают меньшие потери.
    \item \textbf{Связь с нейронауками}: метод разреженного кодирования восходит к
          классической модели Olshausen \& Field~\cite{olshausen1997sparse} представления
          зрительной коры V1 через разреженный базис, что придаёт подходу дополнительную
          биологическую мотивацию.
\end{itemize}

\subsection{Выводы по методу}

Работа «Towards Monosemanticity» вносит существенный вклад в механистическую
интерпретируемость, демонстрируя:
\begin{enumerate}
    \item Разреженные автокодировщики позволяют декомпозировать полисемантические
          нейронные активации в более интерпретируемое пространство моносемантических признаков;
    \item Функция потерь~\eqref{eq:loss} обеспечивает баланс между точностью реконструкции
          и разреженностью при контролируемом $\lambda$;
    \item Выученные признаки каузально влияют на поведение модели (feature steering);
    \item Увеличение размера словаря приводит к дроблению признаков и росту
          интерпретируемости;
    \item Автоматическая оценка интерпретируемости с помощью языковых моделей
          позволяет масштабировать анализ на тысячи признаков.
\end{enumerate}

% ------------------- ЗАКЛЮЧЕНИЕ -------------------
\chapter*{Заключение}
\addcontentsline{toc}{chapter}{Заключение}
Краткий обзор результатов работы:
\begin{itemize}
    \item Достигнута цель...
    \item Решены задачи...
\end{itemize}
\textbf{Перспективы дальнейшей деятельности:} ...

% ------------------- БИБЛИОГРАФИЧЕСКИЙ СПИСОК -------------------
\chapter*{Библиографический список}
\addcontentsline{toc}{chapter}{Библиографический список}
\printbibliography[heading=none]

% ------------------- ПРИЛОЖЕНИЯ -------------------
\appendix
\chapter{Глоссарий}
\section*{Терминологический словарь}
\begin{description}
    \item[Overfitting] Переобучение модели.
    \item[Regularization] Регуляризация.
\end{description}

\chapter{Исходные данные экспериментов}
\section{Протоколы}
Текст протоколов.

\chapter{Листинги кода (фрагменты)}
\begin{verbatim}
import numpy as np
def loss(y_true, y_pred):
    return np.mean((y_true - y_pred)**2)
\end{verbatim}

\end{document}